\section{R3: Repair Quality}
\label{sec:eval-r3}

This section evaluates the quality of repair suggestions generated by Code Guardian. The assessment combines a per-sample manual review on a stratified subset and a fully-automated auto-applicable rate computed on the entire stage-2 output. The manual review covers 25 samples from the best configuration (\texttt{qwen3:8b+RAG}); samples were drawn to span the highest-frequency canonical categories in the curated corpus, with each major category (SQL injection, command injection, path traversal, XSS, SSRF, prototype pollution, race condition) contributing 2--5 samples (Table~\ref{tab:r3-by-type}). The auto-applicable rate is reported separately in Section~\ref{sec:eval-r3-auto-applicable} on the full 279-repair output. R3 thresholds are defined in Section~\ref{sec:r3-repair-quality}.

Of the 25 reviewed samples, 19 (76\%) included a repair suggestion. Among those 19 repairs, 17 (89.5\%) were semantically correct, 18 (94.7\%) aligned with the expected fix strategy, and 8 (42.1\%) contained directly executable code; the remainder provided prose guidance describing the fix approach.

\subsection{Repair Quality by Vulnerability Type}

Repair effectiveness varies by vulnerability category. Table~\ref{tab:r3-by-type} shows fix rates for the most common categories.

\begin{table}[H]
\centering
\caption{Repair generation rate by vulnerability type (\texttt{qwen3:8b+RAG}).}
\label{tab:r3-by-type}
\footnotesize
\begin{tabular}{@{}lrr@{}}
\toprule
\textbf{Vulnerability Type} & \textbf{Samples} & \textbf{Fix Rate (\%)} \\
\midrule
SQL Injection & 5 & 100 \\
Command Injection & 3 & 100 \\
Path Traversal & 3 & 100 \\
XSS & 4 & 75 \\
SSRF & 2 & 50 \\
Prototype Pollution & 2 & 0 \\
Race Condition & 2 & 0 \\
\bottomrule
\end{tabular}
\end{table}

Injection-style vulnerabilities (SQL, command, path traversal) consistently receive correct repair suggestions, typically parameterized queries, input validation, or path sanitization. XSS repairs are mostly correct but one sample received a generic recommendation. Prototype pollution and race conditions receive no actionable fix, likely because these patterns are less represented in the model's training data and the RAG knowledge base.

\subsection{Representative Examples}

\paragraph{True positive with correct repair (SQL injection).}
The system correctly identified an unsanitized \texttt{req.query} value concatenated into a SQL string. The suggested fix replaced string concatenation with a parameterized query using \texttt{db.query("SELECT * FROM users WHERE id = ?", [req.query.id])}, which directly addresses the root cause.

\paragraph{False positive with misleading repair.}
A secure \texttt{execFile} call with a hardcoded command was flagged as a command injection risk. The repair suggested replacing it with \texttt{spawn} and input validation, which is unnecessary since the original code does not use user-controlled input.

\subsection{Level Mapping}

Based on the R3 evaluation scale (Table~\ref{tab:r3-repair-thresholds}):

\begin{itemize}
  \item 76\% of samples received a fix.
  \item Of those, 89.5\% were semantically correct.
  \item Combined: roughly 68\% of all samples received a correct, relevant repair.
\end{itemize}

This falls in the $[50\%, 75\%)$ range.

\begin{center}
\begin{tabular}{cc}
\raisebox{-0.3cm}{\tikz{\filldraw[fill=black] (0,0) -- (90:0.35cm) arc (90:-90:0.35cm) -- cycle; \draw (0,0) circle (0.35cm);}} &
\textbf{R3 Level: Medium repair quality.} \\
\end{tabular}
\end{center}

\noindent\textit{Note on visual scale: R3 uses a three-level scale (High / Medium / Low), as does R4; the half-filled circle above denotes Medium in both. R1 and R2 use a four-level scale (High / Medium / Low / Insufficient) in which the half-filled symbol falls between Medium and Low and is not a formal band --- it appears only in the positioning comparison table (Section~\ref{sec:related-work}), which uses five levels for cross-tool comparison. The per-requirement evaluation sections each define their own scale in the accompanying threshold table.}

Most fixes target the right issue but may use generic patterns or need minor edits before applying.

\paragraph{Limitations of R3 assessment.} All repair assessments were performed by a single reviewer on 25 samples. Inter-rater reliability was not measured, which limits the generalizability of the quality ratings. A two-rater protocol with Cohen's $\kappa$ would strengthen R3 evidence in future evaluations. Additionally, the sample size is small relative to the full dataset; expanding the reviewed set would improve confidence in per-category fix rates.

\subsection{Auto-Applicable Rate}
\label{sec:eval-r3-auto-applicable}

The 25-sample manual review above operates on a generous bar: a repair is scored ``semantically correct'' if the prose or code adequately describes the right fix strategy. A stricter, fully-automated bar is whether the \texttt{suggestedFix} string is itself syntactically applicable code --- i.e.\ whether a developer could accept the IDE quick fix without rewriting the suggestion. The harness reports this as the \emph{auto-applicable rate} (Section~\ref{sec:eval-metrics}).

Repairs are produced by a dedicated stage-2 call per detected finding (mirroring the production extension's split between detection and repair generation in \texttt{src/analyzer.ts}). The repair call uses Ollama's structured-output schema with shape \texttt{\{code: string, language: 'javascript' | 'typescript'\}}, and the system prompt instructs the model to ``Return only executable code in \texttt{code}; do NOT include prose, explanation, or markdown fences.'' The schema therefore constrains the field shape rather than just its type, so prose paraphrases of the fix cannot satisfy the contract.

Across the full \texttt{qwen3:8b+RAG} evaluation (279 stage-2 repairs from 213 vulnerable evaluations under three-pass consensus and the inter-run confidence gate at $\geq 0.67$), \textbf{252 (90.32\,\%) auto-applied} cleanly through the syntax validator. The 27 remaining rejections are all real fragments of code that fail to parse for a structural reason: mid-statement truncations by \texttt{num\_predict} (e.g.\ \texttt{connection.query(sql, [user, pass], (error, results) => \{} with no closing brace) and unterminated regular-expression literals. No rejections come from prose-language openings.

\paragraph{Interpretation.}
The auto-applicable rate measures whether the IDE quick-fix surface could insert each repair without further editing --- a deterministic alternative to the similarity-based metrics (cosine, BLEU, CodeBERTScore) that SecRepair~\cite[\S III-C3]{islam2024secrepair} identifies as ``fundamentally inadequate'' for security correctness. The 9.68\,\% residual is a model-emit defect (truncation, unterminated literal), not a prompting defect. The thesis reports both the manual review (\emph{decision-support quality}: prose fixes that describe the right strategy count) and the auto-applicable rate (\emph{automation quality}: only IDE-insertable fixes count), so a reader can pick the bar that matches their workflow.
